\section{Semantic Routing}

Semantic routing determines which computational pathway a fused and stabilized
semantic state should follow. In the MEM8-Geometric-Architecture, routing is
performed using geometric similarity between the current semantic vector and a
set of learned route vectors. This mechanism enables context-sensitive
transitions across the cognitive manifold.

Let $\mathbf{z} \in \mathbb{R}^d$ denote the fused semantic vector after
stability filtering. Let $\{\mathbf{r}_i\}_{i=1}^k$ be a set of learned route
vectors, each representing a distinct processing pathway or cognitive
operation.

\subsection{Cosine Similarity Routing}

Routing is defined as selecting the route vector with maximum cosine similarity
to the current semantic state:

\[
i^\* = \arg\max_i\,
\cos(\mathbf{z}, \mathbf{r}_i)
\]

where cosine similarity is:

\[
\cos(\mathbf{z}, \mathbf{r}_i) =
\frac{\mathbf{z}^\top \mathbf{r}_i}
{\|\mathbf{z}\|\,\|\mathbf{r}_i\|}
\]

This metric evaluates directional alignment on the semantic manifold, ensuring
that routing decisions reflect geometric coherence rather than magnitude alone.

\subsection{Geometric Interpretation}

Each route vector $\mathbf{r}_i$ defines a direction on the semantic manifold.
Routing selects the direction most aligned with the current semantic state,
interpreting cognition as movement along preferred geometric trajectories.

This formulation supports:

\begin{itemize}
\item curvature-aware transitions,
\item drift-sensitive routing,
\item stability-preserving updates,
\item smooth geodesic continuation.
\end{itemize}

Because cosine similarity is invariant to scale, routing decisions depend only
on direction, making them robust to magnitude fluctuations introduced by
modality imbalance or partial observability.

\subsection{Route Vector Learning}

Route vectors are learned through gradient-based optimization, typically using
contrastive or alignment-based objectives. Each route vector corresponds to a
distinct cognitive function, such as:

\begin{itemize}
\item memory update,
\item contextual expansion,
\item modality reweighting,
\item stability reinforcement,
\item action selection.
\end{itemize}

The geometric separation of route vectors ensures that routing decisions remain
interpretable and stable across diverse semantic inputs.

\subsection{Role in the Cognitive Loop}

Routing occurs after stability filtering and before membrane evaluation. The
selected route determines which geometric operators will be applied next,
including curvature checks, collapse-risk evaluation, Jacobian smoothness
analysis, and firewall enforcement.

This makes routing a critical decision point in the cognitive loop, guiding the
semantic state along appropriate geometric pathways.
